\scanpage{51}
\margin{50}
\mspage{5}
\deleted{
\[
(20)\qquad\uline{\omega}\simeq\{\pm1\}\quad\text{i.e. on dispose }\uline{\omega}\xrightarrow{\sim}\uline{\varepsilon}
\]
donc dorénavant on identifie l'ens $\uline{\omega}=\Omega(J)$ \inserted{$=\Omega(\Sigma_{\mathbb R})$} avec l'ens $\uline{\varepsilon}$ des deux hémisphères.
}
\margin{Il s'agit d'un \unsure{cercle}, \unsure{qui est déterminé|qui est déjà muni} \unsure{origine}, extrémité, \unsure{désignation} du \unsure{parcours à priori}, \unsure{du sens|sur deux} \unsure{points}.}
\margin{\uline{NB}
\[
\operatorname{Rep}(J)\simeq J\times\uline{\omega}
\]
\[
(i,j,k)\longmapsto(i,\omega_{i\to j\to k})\qquad(20)
\]
le paquet $(a_i^{\unsure{\omega\varepsilon}})$ est un torseur sous $\Gamma=\mathcal G_J\times\{1,\tau\}$.
}
Introduisons sur $\widetilde\Sigma$ les \deleted{\unsure{arcs}} chemins simples : (segments orientés plongés)
\[
(21)\qquad\overbrace{a_i^{\omega,\varepsilon}}^{12},\quad
\overbrace{b_i^\omega}^{6},\quad\overbrace{c_i^\varepsilon}^{6},\quad\overbrace{d_i^\varepsilon}^{6},\qquad
i\in J,\ \omega\in\uline{\omega},\ \varepsilon\in\uline{\varepsilon}
\]
chaque paquet \inserted{$b,c,d$} formant un torseur sous $\mathcal G_J$

$\tau$ opérant par
\[
(22)\quad \tau a_i^{\omega,\varepsilon}=a_i^{\omega,\varepsilon'},\quad
\tau b_i^\omega=b_i^\omega,\quad \tau c_i^\varepsilon=c_i^{\varepsilon'},\quad \tau d_i^\varepsilon=d_i^{\varepsilon'}.
\]
\margin{\uline{NB} on a donc \deleted{\unsure{déduit}}
$g(a_i^\omega)=a_{g(i)}^{g(\omega)}$ etc.}
On a
\[
(23)\qquad a_i^{\omega\varepsilon}:S_i^\varepsilon\longrightarrow\deleted{\unsure{S_i^\varepsilon}}R_i^\omega\text{ sur }U_i\quad\inserted{(\text{quart de tour})}
\]
\deleted{dans le sens de l'orientation de $\Sigma$ en $R_i$}
ce qui caractérise $a_i^{\omega\varepsilon}$. \deleted{\unsure{explicite le fait que}}
\[
(24)\qquad b_i^\omega:Q_i\longrightarrow R_{\omega i}^{\omega'}\text{ sur }\widetilde\Sigma_{\mathbb R}\simeq\deleted{\unsure{\Sigma_{\mathbb R}}}\inserted{\operatorname{\text{Déc}}(\Sigma_{\mathbb R},J)}\ \deleted{\unsure{\text{sur les}}}
\]
\deleted{(chemin d'origine \unsure{sur ce cercle})}
(chemin sur $\widetilde\Sigma_{\mathbb R}$, d'origine $Q_i$, dans la direction $\omega$, jusqu'à l'extrémité \deleted{\unsure{qui est}} du segment composante de $\widetilde\Sigma_{\mathbb R}=\operatorname{\text{Déc}}(\Sigma_{\mathbb R},J)$, savoir $R_{\omega i}^{\omega'}$.
\[
(25)\qquad c_i^\omega:P^\omega\mathrel{-\!\!\!-\!\!\!-}Q_i\qquad\text{suivant un arc de grand cercle}
\]
\[
(26)\qquad d_i^\omega:P^\omega\longrightarrow S_i^\omega\qquad\text{suivant un arc de grand cercle}
\]
\margin{On pose aussi
\[
(27)\quad\widetilde b_i^\omega:R_{\omega i}^{\unsure{\omega|\omega'}}\longrightarrow R_{\omega'i}^{\omega'},\quad\widetilde b_i^\omega=b_i^\omega(b_i^{\omega'})^{-1}
\]
chemins dans la direction $\omega$ \unsure{et non pour} $\unsure{b_i}$\deleted{\unsure{des points}}\deleted{\illegible{3}}
}
Tous les arcs $\widetilde b_i^\omega,c_i^\omega,d_i^\omega$ se font sur des grands cercles de $\Sigma$, sauf que les extrémités, des $c_i^\omega,d_i^\omega$ ne sont pas dans $\Sigma^*$ mais dans \inserted{$U_i\subset$} $\widetilde\Sigma\setminus\Sigma^*$
